\PosterSection{4. RESULTADOS E DISCUSSÃO}
\PosterBody
Random Forest obteve o menor MAE macro ($0{,}409 \pm 0{,}062$); LSTM, o menor RMSE macro ($0{,}548 \pm 0{,}073$). A SNN alcançou MAE $0{,}447 \pm 0{,}064$, RMSE $0{,}566 \pm 0{,}072$ e $R^2$ $0{,}291 \pm 0{,}146$: desempenho competitivo para implantação na borda, sem superioridade preditiva demonstrada.

A turbidez apresenta o maior MAE absoluto por sua faixa numérica mais ampla. O OD é o alvo mais difícil em variância explicada: todos os modelos apresentaram $R^2$ médio negativo, coerente com sua baixa correlação com as entradas disponíveis. O pH apresentou o comportamento preditivo mais confiável. Os intervalos resumem a variabilidade entre dez sementes e não representam incerteza hidrológica populacional completa.

\begin{center}
% Vector charts built directly from the multiseed summary and the article's
% platform benchmark table. Values are intentionally explicit for review.
\makebox[\linewidth][c]{%
\begin{tikzpicture}
\begin{axis}[
  width=0.96\linewidth,
  height=8.2cm,
  ybar,
  bar width=15pt,
  ymin=0, ymax=1.15,
  ylabel={MAE macro (menor é melhor)},
  symbolic x coords={RF,LSTM,SNN,MLP},
  xtick=data,
  tick label style={font=\fontsize{28}{32}\selectfont},
  label style={font=\fontsize{30}{34}\selectfont},
  nodes near coords,
  every node near coord/.append style={font=\fontsize{25}{29}\selectfont,anchor=south},
  error bars/y dir=both,
  error bars/y explicit,
  enlarge x limits=0.18,
  axis line style={black!70},
  ymajorgrids=true,
  grid style={gray!25},
  title={\textbf{Desempenho preditivo entre modelos}},
  title style={font=\fontsize{30}{35}\selectfont}
]
\addplot+[draw=ChartGreen,fill=ChartGreen!82] coordinates {
  (RF,0.409) +- (0,0.062)
  (LSTM,0.428) +- (0,0.059)
  (SNN,0.447) +- (0,0.064)
  (MLP,0.781) +- (0,0.193)
};
\end{axis}
\end{tikzpicture}
}

\vspace{0.45cm}
\makebox[\linewidth][c]{%
\begin{tikzpicture}
\begin{axis}[
  width=0.96\linewidth,
  height=8.5cm,
  ymode=log,
  ymin=0.001, ymax=20,
  ylabel={Energia incremental (mJ/inferência)},
  symbolic x coords={ESP32,RPi4,AndA,AndB,AndC},
  xtick=data,
  xticklabels={ESP32,Raspberry\ Pi 4B,Android\ A,Android\ B,Android\ C},
  tick label style={font=\fontsize{25}{29}\selectfont,align=center},
  label style={font=\fontsize{30}{34}\selectfont},
  point meta=explicit symbolic,
  nodes near coords,
  every node near coord/.append style={font=\fontsize{25}{29}\selectfont,anchor=west,xshift=6pt},
  enlarge x limits=0.14,
  axis line style={black!70},
  ymajorgrids=true,
  grid style={gray!25},
  title={\textbf{Energia por inferência nas plataformas}},
  title style={font=\fontsize{30}{35}\selectfont}
]
\addplot+[only marks,mark=*,mark size=5.5pt,draw=ChartBlue,fill=ChartBlue] coordinates {
  (ESP32,11.204) [11{,}204]
  (RPi4,0.0512) [0{,}0512]
  (AndA,0.0104) [0{,}0104]
  (AndB,0.0151) [0{,}0151]
  (AndC,0.00285) [0{,}00285]
};
\end{axis}
\end{tikzpicture}
}

\end{center}
\FigureSource{MAE macro (desvio-padrão, dez sementes) e energia incremental por inferência.}
